\chapter{Fragile X Syndrome}

\section*{Definition and Overview}
Fragile X syndrome is the most common inherited cause of intellectual disability and the most common known single-gene cause of autism spectrum disorder. It is caused by a change in the FMR1 gene on the X chromosome, which normally produces a protein essential for brain development. Affected people have intellectual disability of varying severity, characteristic physical features, and behavioural difficulties including anxiety, hyperactivity, and autistic traits. The condition affects males more severely than females because it is X-linked. There is no cure, but early diagnosis, educational support, and treatment of associated problems greatly improve outcomes. This chapter explains the genetics, features, and management of fragile X syndrome.

\section*{Epidemiology}
Fragile X syndrome affects around 1 in 4,000 males and 1 in 8,000 females worldwide. It is the most common inherited form of intellectual disability and accounts for a small but important proportion of autism diagnoses. Around 1 in 250 women carries a premutation of the FMR1 gene, which can expand when passed to children, causing the full syndrome; premutation carriers can also be affected by related conditions, including fragile X-associated tremor/ataxia syndrome (FXTAS) in older adults and premature ovarian insufficiency in women. Many cases are still diagnosed late or missed, despite the availability of reliable genetic testing.

\section*{Aetiology and Pathophysiology}
\begin{itemize}
    \item \textbf{The FMR1 gene:} on the X chromosome, it produces the fragile X mental retardation protein (FMRP), which is essential for normal brain development and synaptic function
    \item \textbf{The mutation:} a repeated DNA sequence (CGG repeat) expands abnormally, usually beyond 200 repeats, silencing the gene so FMRP is not produced
    \item \textbf{Premutation:} 55--200 repeats is a premutation, which does not cause fragile X syndrome but can expand to the full mutation when passed from a mother to her children
    \item \textbf{Inheritance:} X-linked; males with the full mutation are usually more severely affected because they have only one X chromosome, while females have a second X chromosome that partially compensates
    \item \textbf{Effects on the brain:} without FMRP, synapses do not develop and prune normally, causing intellectual disability and behavioural features
    \item \textbf{Genetic anticipation:} the repeat can expand further in successive generations, so affected families can show increasing severity
    \item \textbf{Carrier conditions:} premutation carriers can develop FXTAS (tremor and balance problems in later life) and fragile X-associated primary ovarian insufficiency
\end{itemize}

\section*{Clinical Features}
\begin{itemize}
    \item Intellectual disability, ranging from mild learning difficulties to severe disability; more marked in males
    \item Behavioural features: anxiety, hyperactivity, poor attention, hand flapping, poor eye contact, and autistic traits
    \item Speech and language delay, with repetitive speech
    \item Physical features: a long, narrow face, large ears, a prominent jaw, and large testicles after puberty
    \item Loose joints, flat feet, and soft skin
    \item Medical problems: recurrent ear infections, seizures in a minority, and strabismus
    \item In females: milder features, sometimes only learning difficulties or anxiety
    \item Sensory sensitivities and difficulty with transitions and change
\end{itemize}

\section*{Differential Diagnosis}
Other causes of intellectual disability and autism must be considered.
\begin{itemize}
    \item Other genetic causes of intellectual disability and autism, including chromosomal abnormalities
    \item Autism spectrum disorder without a known genetic cause
    \item Attention deficit hyperactivity disorder (ADHD)
    \item Rett syndrome, Down syndrome, and other genetic syndromes
    \item Environmental causes of developmental delay
    \item Fragile X syndrome is confirmed by genetic testing and can be distinguished from all of these
\end{itemize}

\section*{When to Seek Medical Care}
Families should seek assessment if a child has developmental delay, intellectual disability, autism features, or a family history of fragile X syndrome or intellectual disability. Early diagnosis opens access to early intervention and support.

\textbf{Call 000 (or local emergency services) for:}
\begin{itemize}
    \item Seizures or fitting
    \item Any medical emergency in a child or adult with fragile X syndrome
    \item Suicidal thoughts or severe distress, which can occur with associated anxiety and depression
\end{itemize}

\section*{Diagnosis and Investigations}
\begin{itemize}
    \item \textbf{Genetic testing:} a blood test measuring the FMR1 gene, which confirms the full mutation or premutation; this is the definitive diagnostic test
    \item \textbf{Developmental assessment:} cognitive, language, and behavioural assessment by specialists
    \item \textbf{Family testing:} carrier testing of at-risk relatives and prenatal or preimplantation diagnosis for families
    \item \textbf{Associated assessments:} hearing and vision tests, and investigations for seizures and other medical problems
    \item \textbf{Early diagnosis:} allows early intervention, which improves developmental outcomes
\end{itemize}

\section*{Management}
\begin{itemize}
    \item \textbf{Early intervention:} speech therapy, occupational therapy, physiotherapy, and early childhood programs
    \item \textbf{Educational support:} individualised plans with structure, visual supports, and predictable routines
    \item \textbf{Behavioural support:} strategies for anxiety, attention, and sensory sensitivities, with psychology input
    \item \textbf{Medicines:} for co-occurring ADHD, anxiety, and seizures, under medical supervision
    \item \textbf{Speech and communication:} ongoing speech therapy, including augmentative communication where needed
    \item \textbf{Treat medical problems:} ear infections, vision problems, and seizures as they arise
    \item \textbf{Transition and adult support:} vocational training, employment support, and supported living as needed
    \item \textbf{Genetic counselling:} for families, covering recurrence risks and carrier implications
    \item \textbf{NDIS support:} funding for therapies and supports through the National Disability Insurance Scheme
\end{itemize}

\section*{Complications}
\begin{itemize}
    \item Anxiety, depression, and behavioural difficulties, which are common
    \item Epilepsy in a minority
    \item Recurrent infections, including ear infections in childhood
    \item Premature ovarian insufficiency in female premutation carriers
    \item FXTAS (tremor, balance, and cognitive problems) in older premutation carriers
    \item The family burden of caring for a person with significant disability
\end{itemize}

\section*{Prognosis and Outcomes}
Fragile X syndrome is a lifelong condition, but outcomes vary widely and have improved with early diagnosis and support. People with fragile X syndrome can learn, form relationships, and participate in community life, with many achieving literacy, employment, and independent living to varying degrees. Early intervention and structured education improve outcomes substantially, and treatment of anxiety and ADHD improves quality of life. Research into treatments that restore FMRP function is underway. For families, genetic counselling and support networks make a real difference to the journey.

\section*{Prevention and Self-Care}
\begin{itemize}
    \item Seek genetic testing early for developmental delay, intellectual disability, or autism with suggestive features
    \item Offer carrier testing and genetic counselling to at-risk family members
    \item Consider prenatal or preimplantation genetic testing in affected families
    \item Access early intervention services and the NDIS promptly
    \item Build structure, routine, and visual supports into daily life
    \item Support the person's mental health and treat anxiety early
    \item Connect with fragile X support organisations for information and community
    \item Look after carers' wellbeing, including respite and support groups
\end{itemize}

\section*{Sources and Further Reading}
The following reputable sources provide further information for healthcare students and patients seeking reliable, current advice about fragile X syndrome.
\begin{itemize}
    \item Fragile X Association of Australia. \textit{Information and support.} https://www.fragilex.org.au/
    \item Genetic Alliance Australia. \textit{Genetic conditions.} https://www.geneticalliance.org.au/
    \item Healthdirect Australia. \textit{Fragile X syndrome.} https://www.healthdirect.gov.au/
    \item Centre for Genetics Education. \textit{Fragile X syndrome.} https://www.genetics.edu.au/
    \item NHS. \textit{Fragile X syndrome.} https://www.nhs.uk/conditions/fragile-x-syndrome/
    \item MedlinePlus. \textit{Fragile X syndrome.} https://medlineplus.gov/
\end{itemize}
